\documentclass[tikz,border=6pt]{standalone}
\usepackage{amsmath,amssymb,mathtools,physics,bm,nicefrac,wasysym}
\usepackage{xcolor}

\definecolor{colone}{RGB}{180,60,60}
\definecolor{coltwo}{RGB}{60,110,180}
\definecolor{colthree}{RGB}{60,140,80}

\definecolor{accent}{HTML}{A13D2C}
\definecolor{accent2}{HTML}{2B5F6B}
\definecolor{muted}{HTML}{5A564F}
\definecolor{panel}{HTML}{F8F6F0}
\definecolor{linecol}{HTML}{D8D2C4}
\definecolor{ink}{HTML}{1E1C1A}
\definecolor{astroblue}{HTML}{1B4F72}
\definecolor{astrodark}{HTML}{17202A}
\definecolor{sunyellow}{HTML}{E8A33D}

\usetikzlibrary{calc,arrows.meta,positioning,decorations.pathreplacing,decorations.markings,angles,quotes,shapes.geometric}
\usepackage{tikz-3dplot}

\newcommand{\vecr}{\vec r}
\newcommand{\vecL}{\vec L}
\newcommand{\vecA}{\vec A}
\newcommand{\vecf}{\vec f}
\newcommand{\vecF}{\vec F}
\newcommand{\vecp}{\vec p}
\newcommand{\avg}[1]{\left\langle #1\right\rangle}
\newcommand{\vecx}{\vec x}
\newcommand{\hatn}{\hat n}
\newcommand{\rhat}{\hat r}
\newcommand{\phihat}{\hat\phi}
\newcommand{\thetahat}{\hat\theta}
\newcommand{\note}[1]{\textcolor{blue!60!black}{\textit{[#1]}}}

% Astronomical acronyms
\newcommand{\LST}{\mathrm{LST}}
\newcommand{\GST}{\mathrm{GST}}
\newcommand{\GMST}{\mathrm{GMST}}
\newcommand{\GAST}{\mathrm{GAST}}
\newcommand{\LMST}{\mathrm{LMST}}
\newcommand{\LAST}{\mathrm{LAST}}
\newcommand{\UT}{\mathrm{UT}}
\newcommand{\UTC}{\mathrm{UTC}}
\newcommand{\TAI}{\mathrm{TAI}}
\newcommand{\TT}{\mathrm{TT}}
\newcommand{\TDB}{\mathrm{TDB}}
\newcommand{\JD}{\mathrm{JD}}
\newcommand{\MJD}{\mathrm{MJD}}
\newcommand{\EoT}{\mathrm{EoT}}
\newcommand{\SNR}{\mathrm{SNR}}
\newcommand{\RON}{\mathrm{RON}}

% Helper commands for specific diagrams
\newcommand{\midlinelabel}[3]{
    \node (midlabel) at ($ (#1)!.5!(#2) $) {#3};
    \draw[stealth-] (#1) --  (midlabel);
    \draw[-stealth] (midlabel) -- (#2);
}
\newcommand{\midlinelabell}[3]{
    \node[fill = white, inner sep = 0.2pt, outer sep=3pt] (midlabel) at ($ (#1)!.5!(#2) $) {#3};
    \draw[stealth-] (#1) --  (midlabel);
    \draw[-stealth] (midlabel) -- (#2);
}
\newcommand{\midlinelabelll}[3]{
    \node[fill = white, inner sep = 0.2pt, outer sep=3pt,scale=0.6] (midlabel) at ($ (#1)!.5!(#2) $) {#3};
    \draw[stealth-] (#1) --  (midlabel);
    \draw[-stealth] (midlabel) -- (#2);
}

\tdplotsetmaincoords{70}{120}

\begin{document}
\tdplotsetmaincoords{70}{120}
\begin{tikzpicture}[tdplot_main_coords, scale=3]
				
				%------------------------------------------------
				% Parameters
				%------------------------------------------------
				\def\a{2.7}
				\def\e{0.45}
				\def\inc{35}
				\def\Om{35}
				\def\om{20}
				
				% Position of the whole orbit
				\def\Ox{0.0}
				\def\Oy{0.0}
				\def\Oz{0.0}
				
				% Point to mark
				\def\numark{55}
				
				%------------------------------------------------
				% Sky plane
				%------------------------------------------------
				\fill[blue!5]
				(-4.5,-3.8,0) --
				( 4.5,-3.8,0) --
				( 4.5, 3.8,0) --
				(-4.5, 3.8,0) -- cycle;
				
				% X,Y axes: reference / sky plane
				\draw[->,blue!70!black]
				(-4.5,0,0) -- (4.8,0,0)
				node[anchor=north east] {$\hat\delta$};
				
				\draw[->,blue!70!black]
				(0,4.2,0) -- (0,-3.8,0)
				node[anchor=south] {$\hat\alpha$};
				
				% Z axis
				\draw[->,gray!70!black]
				(0,0,-1.5) -- (0,0,3.0)
				node[anchor=south west] {$Z$};
				
				\node[blue!60!black] at (3.2,3.2,0)
				{\small sky plane};
				
				%------------------------------------------------
				% Origin / centre of orbit
				%------------------------------------------------
				
				\fill (\Ox,\Oy,\Oz) circle (1.5pt);
				
				\node[above left] at (\Ox,\Oy,\Oz) {$O$};
				
				%------------------------------------------------
				% Omega
				% Angle from +alpha-hat toward +delta-hat
				%------------------------------------------------
				
				\draw[thick,blue!70!black,->]
				plot[
				variable=\theta,
				domain=-90:\Om,
				samples=40
				]
				(
				{\Ox+0.65*cos(\theta)},
				{\Oy+0.65*sin(\theta)},
				{\Oz}
				);
				
				
				\node[blue!70!black] at
				({\Ox+0.85*cos(\Om/2)},
				{\Oy+0.85*sin(\Om/2)},0)
				{$\Omega$};
				%------------------------------------------------
				% Actual inclined ellipse
				% Orbital plane
				%------------------------------------------------
				
				\draw[thick,red!70!black]
				plot[
				variable=\nu,
				domain=0:360,
				samples=250,
				smooth
				]
				(
				{
					\Ox+
					(\a*(1-\e*\e)/(1+\e*cos(\nu)))
					*
					(
					cos(\Om)*cos(\om+\nu)
					-sin(\Om)*cos(\inc)*sin(\om+\nu)
					)
				},
				{
					\Oy+
					(\a*(1-\e*\e)/(1+\e*cos(\nu)))
					*
					(
					sin(\Om)*cos(\om+\nu)
					+cos(\Om)*cos(\inc)*sin(\om+\nu)
					)
				},
				{
					\Oz+
					(\a*(1-\e*\e)/(1+\e*cos(\nu)))
					*
					(
					sin(\inc)*sin(\om+\nu)
					)
				}
				);
				
				%------------------------------------------------
				% Projected ellipse
				% Lies in XY sky plane
				%------------------------------------------------
				
				\draw[very thick,dashed,blue!45!black]
				plot[
				variable=\nu,
				domain=0:360,
				samples=250,
				smooth
				]
				(
				{
					\Ox+
					(\a*(1-\e*\e)/(1+\e*cos(\nu)))
					*
					(
					cos(\Om)*cos(\om+\nu)
					-sin(\Om)*cos(\inc)*sin(\om+\nu)
					)
				},
				{
					\Oy+
					(\a*(1-\e*\e)/(1+\e*cos(\nu)))
					*
					(
					sin(\Om)*cos(\om+\nu)
					+cos(\Om)*cos(\inc)*sin(\om+\nu)
					)
				},
				0
				);
				
				%------------------------------------------------
				% Periapsis
				%------------------------------------------------
				
				\pgfmathsetmacro{\rp}{\a*(1-\e)}
				
				\pgfmathsetmacro{\periX}
				{\Ox+\rp*(cos(\Om)*cos(\om)
					-sin(\Om)*cos(\inc)*sin(\om))}
				
				\pgfmathsetmacro{\periY}
				{\Oy+\rp*(sin(\Om)*cos(\om)
					+cos(\Om)*cos(\inc)*sin(\om))}
				
				\pgfmathsetmacro{\periZ}
				{\Oz+\rp*sin(\inc)*sin(\om)}
				
				% Direction to periapsis
				\draw[thick,red!70!black,->]
				(\Ox,\Oy,\Oz)
				-- (\periX,\periY,\periZ);
				
				% Periapsis point
				\fill[red!70!black]
				(\periX,\periY,\periZ) circle (2pt);
				
				\node[red!70!black,above right] at
				(\periX,\periY,\periZ)
				{periapsis};
				
				%------------------------------------------------
				% omega
				% Angle in orbital plane
				%------------------------------------------------
				
				\def\Romega{0.75}
				
				\draw[thick,red!70!black,->]
				plot[
				variable=\theta,
				domain=0:\om,
				samples=30
				]
				(
				{
					\Ox+\Romega*
					(
					cos(\Om)*cos(\theta)
					-sin(\Om)*cos(\inc)*sin(\theta)
					)
				},
				{
					\Oy+\Romega*
					(
					sin(\Om)*cos(\theta)
					+cos(\Om)*cos(\inc)*sin(\theta)
					)
				},
				{
					\Oz+\Romega*sin(\inc)*sin(\theta)
				}
				);
				
				\node[red!70!black] at
				(
				{
					\Ox+1.0*
					(
					cos(\Om)*cos(\om/2)
					-sin(\Om)*cos(\inc)*sin(\om/2)
					)
				},
				{
					\Oy+1.0*
					(
					sin(\Om)*cos(\om/2)
					+cos(\Om)*cos(\inc)*sin(\om/2)
					)
				},
				{
					\Oz+1.0*sin(\inc)*sin(\om/2)
				}
				)
				{$\omega$};
				
				%------------------------------------------------
				% Inclination I
				% Angle between reference and orbital planes
				%------------------------------------------------
				
				\def\Rinc{1.5}
				
				% Reference-plane direction
				\coordinate (A) at
				({\Rinc*cos(\Om)},
				{\Rinc*sin(\Om)},0);
				
				% Corresponding direction in orbital plane
				\coordinate (B) at
				({\Rinc*cos(\Om)-1.0*sin(\Om)*cos(\inc)},
				{\Rinc*sin(\Om)+1.0*cos(\Om)*cos(\inc)},
				{1.0*sin(\inc)});
				
				% Direction in reference plane
				\coordinate (C) at
				({\Rinc*cos(\Om)-1.0*sin(\Om)},
				{\Rinc*sin(\Om)+1.0*cos(\Om)},0);
				
				% Show the two planes locally
				\draw[blue!50!black,dashed] (A) -- (C);
				\draw[red!50!black,dashed] (A) -- (B);
				
				% Arc representing inclination
				\draw[thick,red!70!black,->]
				plot[
				variable=\theta,
				domain=0:\inc,
				samples=30
				]
				(
				{\Rinc*cos(\Om)-0.55*sin(\Om)*cos(\theta)},
				{\Rinc*sin(\Om)+0.55*cos(\Om)*cos(\theta)},
				{0.55*sin(\theta)}
				);
				
				\node[red!70!black] at
				(
				{\Rinc*cos(\Om)-0.72*sin(\Om)*cos(\inc/2)},
				{\Rinc*sin(\Om)+0.72*cos(\Om)*cos(\inc/2)},
				{0.72*sin(\inc/2)}
				)
				{$I$};
				
				%------------------------------------------------
				% Mark a point P on the actual orbit
				%------------------------------------------------
				
				\pgfmathsetmacro{\rr}
				{\a*(1-\e*\e)/(1+\e*cos(\numark))}
				
				\pgfmathsetmacro{\PX}
				{\Ox+\rr*(cos(\Om)*cos(\om+\numark)
					-sin(\Om)*cos(\inc)*sin(\om+\numark))}
				
				\pgfmathsetmacro{\PY}
				{\Oy+\rr*(sin(\Om)*cos(\om+\numark)
					+cos(\Om)*cos(\inc)*sin(\om+\numark))}
				
				\pgfmathsetmacro{\PZ}
				{\Oz+\rr*sin(\inc)*sin(\om+\numark)}
				
				% Point on orbit
				\fill[red!70!black]
				(\PX,\PY,\PZ) circle (2pt);
				
				\node[red!70!black,above right] at
				(\PX,\PY,\PZ) {$P$};
				
				%------------------------------------------------
				% Projection of P onto sky plane
				%------------------------------------------------
				
				\fill[blue!70!black]
				(\PX,\PY,0) circle (2pt);
				
				\node[blue!70!black,below right] at
				(\PX,\PY,0) {$P_{\rm sky}$};
				
				% Projection line
				\draw[gray,dotted,thick]
				(\PX,\PY,\PZ) -- (\PX,\PY,0);
				
				%------------------------------------------------
				% Direction from O to P
				%------------------------------------------------
				
				\draw[red!40!black,dashed]
				(\Ox,\Oy,\Oz) -- (\PX,\PY,\PZ);
				
				%------------------------------------------------
				% nu
				% Angle from periapsis to P
				% in orbital plane
				%------------------------------------------------
				
				\def\Rnu{1.05}
				
				\draw[thick,red!70!black,->]
				plot[
				variable=\theta,
				domain=\om:\om+\numark,
				samples=50,
				smooth
				]
				(
				{
					\Ox+\Rnu*
					(
					cos(\Om)*cos(\theta)
					-sin(\Om)*cos(\inc)*sin(\theta)
					)
				},
				{
					\Oy+\Rnu*
					(
					sin(\Om)*cos(\theta)
					+cos(\Om)*cos(\inc)*sin(\theta)
					)
				},
				{
					\Oz+\Rnu*sin(\inc)*sin(\theta)
				}
				);
				
				\node[red!70!black] at
				(
				{
					\Ox+1.25*
					(
					cos(\Om)*cos(\om+\numark/2)
					-sin(\Om)*cos(\inc)*sin(\om+\numark/2)
					)
				},
				{
					\Oy+1.25*
					(
					sin(\Om)*cos(\om+\numark/2)
					+cos(\Om)*cos(\inc)*sin(\om+\numark/2)
					)
				},
				{
					\Oz+1.25*sin(\inc)*sin(\om+\numark/2)
				}
				)
				{$\nu$};
				%------------------------------------------------
				% Orbital-plane basis vectors
				%------------------------------------------------
				
				\def\Ry{2.0}
				\def\RL{2.4}
				
				%------------------------------------------------
				% y-hat
				%
				% y-hat =
				% -sin(Omega) cos(i) X
				% +cos(Omega) cos(i) Y
				% +sin(i) Z
				%------------------------------------------------
				
				\pgfmathsetmacro{\yX}
				{\Ox-\Ry*sin(\Om)*cos(\inc)}
				
				\pgfmathsetmacro{\yY}
				{\Oy+\Ry*cos(\Om)*cos(\inc)}
				
				\pgfmathsetmacro{\yZ}
				{\Oz+\Ry*sin(\inc)}
				
				\coordinate (Yhat) at (\yX,\yY,\yZ);
				
				\draw[
				very thick,
				purple!80!black,
				->,
				>=stealth
				]
				(\Ox,\Oy,\Oz) -- (Yhat)
				node[above right] {$\hat y$};
				
				
				%------------------------------------------------
				% L-hat
				%
				% L-hat = n-hat x y-hat
				%
				% L-hat =
				% +sin(Omega) sin(i) X
				% -cos(Omega) sin(i) Y
				% +cos(i) Z
				%------------------------------------------------
				
				\pgfmathsetmacro{\LX}
				{\Ox+\RL*sin(\Om)*sin(\inc)}
				
				\pgfmathsetmacro{\LY}
				{\Oy-\RL*cos(\Om)*sin(\inc)}
				
				\pgfmathsetmacro{\LZ}
				{\Oz+\RL*cos(\inc)}
				
				\coordinate (Lhat) at (\LX,\LY,\LZ);
				
				\draw[
				very thick,
				red!70!black,
				->,
				>=stealth
				]
				(\Ox,\Oy,\Oz) -- (Lhat)
				node[above right] {$\hat L$};
				%------------------------------------------------
				% Line of ascending node
				% Reference-plane quantity
				%------------------------------------------------
				
				\draw[very thick,red!70!black,->](\Ox,\Oy,\Oz)-- ({\Ox+2.4*cos(\Om)},{\Oy+2.4*sin(\Om)},\Oz)node[right] {$\hatn$};
				
		\end{tikzpicture}
\end{document}
